\section{Initial pilot computation accounting}
\label{app:pilot}

Figure~\ref{fig:pilot-budget-accounting} and Table~\ref{tab:pilot-call-costs} document the initial implementation pilot. They expose the distinction between requested coalition-work settings, actual model-input rows and measured call cost. They are not the main reliability evaluation. The source records, panel/cell mappings and unchanged executable presentation generator accompany the artifacts; these calls must remain separate from the main-study inputs and repetitions.

The pilot's smaller games reached exhaustive Kernel SHAP evaluation at the higher requested setting. Complete coverage of a permutation estimator's coalition cache has a different meaning: random path frequencies can still vary. These observations motivate recording actual work and effective-player exhaustion rather than equating a requested setting with a stochastic sample size or a fixed latency. They do not establish an equal-cost accuracy advantage for either estimator.

These initial timings precede the analysis owner's correction to use float64 prediction inputs consistently across all estimators and references. Subsequent timing and full-study accounting must use that corrected implementation; the raw times here describe the archived pilot only. The initial prototype's reliability metrics are not used to estimate the main operational endpoint.

\begin{figure}[htbp]
\centering
\includegraphics[width=\linewidth]{figures/pilot_budget_accounting.pdf}
\caption[Pilot requested and actual prediction work.]{Pilot budget accounting, not main-study reliability evidence. Each line joins two measured calls for one of two pilot inputs per fitted model, with a fixed 32-row empirical background and grouped raw-feature players. The horizontal axis is the requested coalition-work setting; the vertical axis counts actual input rows evaluated by the prediction pipeline. Colors, markers and line styles identify the estimators. Hollow Kernel SHAP markers denote exhaustive evaluation over the effective varying players; complete permutation-cache coverage does not imply an exact path-sampling estimate. Overlapping calls may obscure one another. Traces for the three fitted models coincide in prediction-row count; the accompanying table shows their distinct CPU costs. These pilot calls have no population uncertainty interval. They motivate measuring actual work and using dimension-specific stochastic budgets in the subsequent study.
}
\label{fig:pilot-budget-accounting}
\end{figure}
\clearpage
\input{tables/pilot_call_costs.tex}
\clearpage
